\ssr{ЗАКЛЮЧЕНИЕ}

В процессе выполнения лабораторной работы поставленная цель была достигнута: проведён сравнительный анализ стандартного алгоритма умножения матриц, базового алгоритма Винограда и его оптимизированной модификации по теоретической трудоёмкости и процессорному времени выполнения.

Все задачи работы решены:
\begin{itemize}
    \item описана решаемая задача умножения матриц согласно варианту;
    \item разработаны стандартный алгоритм умножения матриц, алгоритм Винограда и оптимизированный по варианту алгоритм Винограда;
    \item описаны разработанные алгоритмы схемами алгоритмов;
    \item сформулированы требования к разрабатываемому программному обеспечению;
    \item описаны средства реализации алгоритмов;
    \item программно реализованы все разработанные алгоритмы умножения матриц;
    \item проведено функциональное тестирование реализаций;
    \item сформирована модель вычислений и выполнен аналитический расчёт трудоёмкости всех исследуемых алгоритмов;
    \item описаны инструменты замера процессорного времени;
    \item выполнены серии массированных замеров процессорного времени выполнения реализаций для квадратных матриц чётных и нечётных порядков;
    \item проведён сравнительный анализ полученных теоретических оценок трудоёмкости и экспериментальных временных характеристик.
\end{itemize}

В результате исследования установлено, что все рассмотренные алгоритмы обладают кубической асимптотической сложностью $O(N^3)$. В соответствии с теоретической моделью вычислений базовый алгоритм Винограда характеризуется большей трудоёмкостью ($14{,}5N^3$) по сравнению со стандартным ($11N^3$) по старшему члену, однако применение оптимизаций снижает её до $10{,}5N^3$. Экспериментальные замеры подтвердили теоретический порядок роста $O(N^3)$ и отсутствие качественных различий между чётными и нечётными порядками матриц.